% concatenated from Larged2.tex
\documentclass[a4paper,12pt]{article}
\usepackage{amsmath, amsthm, amstext}
\usepackage{amssymb, array, amsfonts}
\usepackage[utf8]{inputenc}
\usepackage{graphicx}

\numberwithin{equation}{section}

\setlength{\topmargin}{-1cm} \setlength{\textheight}{24cm}
\setlength{\oddsidemargin}{-0.5cm}
\setlength{\evensidemargin}{-0.5cm} \setlength{\textwidth}{17cm}

\def \al{\alpha}
\def \be{\beta}
\def \ga{\gamma}
\def \de{\delta}
\def \er{\varepsilon}
\def \ze{\zeta}
\def \ka{\varkappa}
\def \la{\lambda}
\def \si{\sigma}
\def \te{\theta}
\def \ph{\varphi}
\def \oo{\omega}
\renewcommand{\l}{\left}
\renewcommand{\r}{\right}

\def \G{\Gamma}
\def \D{\Delta}
\def \L{\Lambda}
\def \O{\Omega}

\def \C{\mathbb{C}}
\def \N{\mathbb{N}}
\def \R{\mathbb{R}}
\def \Z{\mathbb{Z}}
\def \T{\mathbb{T}}
\def\n{\nabla}
\def\dd{\partial}
\def\div{\operatorname{div}}
\def\rot{\operatorname{rot}}
\def\1{1\!\!\!\!1}
\def\const{\operatorname{const}}
\def\det{\operatorname{det}}
\def\dom{\operatorname{Dom}}
\def\hol{\operatorname{Hol}}
\def\lip{\operatorname{Lip}}
\def\mat{\operatorname{Mat}}
\def\mes{\operatorname{mes}}
\def\ran{\operatorname{Ran}}
\def\spec{\operatorname{spec}}
\def\supp{\operatorname{supp}}
\def\tr{\operatorname{tr}}
\def\im{\operatorname{Im}}
\def\re{\operatorname{Re}}
\newcommand{\<}{\langle}
\renewcommand{\>}{\rangle}

\theoremstyle{plain}
\newtheorem{theorem}{\bf Theorem}[section]
\newtheorem{lemma}[theorem]{\bf Lemma}

\newtheorem{cor}[theorem]{\bf Corollary}

\theoremstyle{definition}
\newtheorem{defi}[theorem]{\bf Definition}

\theoremstyle{remark}
\newtheorem{rem}[theorem]{\bf Remark}

\newtheorem{conj}[theorem]{\bf Conjecture}

\renewcommand{\le}{\leqslant}
\renewcommand{\ge}{\geqslant}

%\renewcommand{\proofname}{Proof}
\renewcommand{\qed}{\vrule height7pt width5pt depth0pt}


\title{Inequalities between Dirichlet and Neumann eigenvalues in large dimensions}
\author{N.~D.~Filonov}
\date{}
\begin{document}
\maketitle

\begin{abstract}
Let $\O$ be a bounded domain in $\R^d$.
Denote by $\la_k$ (resp. $\mu_k$)
the eigenvalues of the Laplace operator in $\O$ with Dirichlet (resp. Neumann)
boundary conditions. 
Denote by $\Psi = \Psi (d,k,\O)$ the shift of indices in the inequality 
$\mu_{k+\Psi} \le \la_k$.
We are interested to describe the behaviour of $\Psi$ for large $d$.
We prove that
a) $\Psi (d,1,\O) \ge C (e/2)^d$ for all domains $\O$; and
b) $\Psi (d,k,\O) \ge C (e/2)^d$ for all $k$ and all convex domains $\O$.
\footnote{Keywords: 
Laplace operator, Dirichlet problem, Neumann problem,
inequalities for eigenavlues, convex domains}
\end{abstract}

\section{Introduction}
Let $d \ge 2$, $\O \subset \R^d$ be a bounded domain 
such that the embedding $W_2^1 (\O) \subset L_2 (\O)$ is compact,
where $W_2^1 (\O)$ is the Sobolev space.
We consider Dirichlet and Neumann problems for the Laplace operator in $\O$:
$$
\begin{cases}
- \D \ph_k = \la_k \ph_k, \\
\left.\ph_k\right|_{\dd\O} = 0,
\end{cases}
\qquad
\begin{cases}
- \D \psi_k = \mu_k \psi_k, \\
\left.\frac{\dd\psi_k}{\dd\nu}\right|_{\dd\O} = 0.
\end{cases}
$$
Spectra of these two problems are discrete,
the eigenvalues form two sequences
\begin{equation*}
0 < \la_1 < \la_2 \le \la_3 \le \dots, \qquad \la_k \to + \infty, 
\end{equation*}
\begin{equation*}
0 = \mu_1 < \mu_2 \le \mu_3 \le \dots, \qquad \mu_k \to + \infty.
\end{equation*}
We take multiplicity into account.
It is clear from the min-max principle that $\mu_k \le \la_k$ always.
Many papers are devoted to the improvements of this inequality,
see \cite{P52, Payne55, Av86, LW86, Fr91, F04, R25}.
In particular, the inequality
\begin{equation}
\label{11}
\mu_{k+d} \le \la_k \qquad \forall \ k \in \N, \ 
\text{for all convex domains} \ \O
\end{equation}
is proved in \cite{LW86}.
The last sentence of this paper says

{\it perhaps \eqref{11} can be replaced by a better inequality of the form
$$
\mu_{\Phi(d,k)} < \la_k
$$
for convex $d$-dimensional domains.}

We show that in large dimensions this inequality indeed can be improved,
see Theorem \ref{t14} below.

Introduce notations
$$
\Phi (d, k, \O) := \# \left\{ j : \mu_j (\O) \le \la_k (\O)\right\},
$$
$$
\Psi (d, k, \O) := \Phi (d, k, \O) - k
$$
in such a way that the inequality
$$
\mu_{k + \Psi(d,k,\O)} \le \la_k
$$
holds true by definition.
In these terms Friedlander's inequality 
$\mu_{k+1} < \la_k$ (see \cite{Fr91, F04}) means 
\begin{equation}
\label{*}
\Psi (d,k,\O) \ge 1
\end{equation}
for all $k$ and all domains.
Levine-Weinberger's inequality means $\Psi (d,k,\O) \ge d$ 
for all $k$ and all {\it convex} domains.

Our first result concerns the cases $k=1$ and $k=2$.

\begin{theorem}
\label{t11}
The inequalities
\begin{equation}
\label{12}
\Phi (d,1,\O) \ge \frac{(j_{d/2-1})^d}{(d+2) 2^{d-1} \G(d/2+1)^2}
\end{equation}
and
\begin{equation*}
\Phi (d,2,\O) \ge \frac{(j_{d/2-1})^d}{(d+2) 2^{d-2} \G(d/2+1)^2}.
\end{equation*}
hold true for all domains $\O$.
Here and everywhere below we denote by $j_\nu$ the first positive root
of the Bessel function $J_\nu$. 
\end{theorem}

\begin{cor}
\label{c12}
There is an absolute constant $C_1 > 0$ such that
\begin{equation}
\label{13} 
\Phi (d,1,\O) \ge C_1 \left(\frac{e}2\right)^d \qquad \forall \ \O.
\end{equation}
\end{cor}

\begin{cor}
\label{c13}
For all domains $\O \in \R^d$ we have
$$
\mu_3 < \la_1 \qquad \text{if} \ \ d \ge 7,
$$
and
$$
\mu_4 < \la_2 \qquad \text{if} \ \ d \ge 6.
$$
\end{cor}

See also the table at the end of the proof of Corollary \ref{c13} below.

Note also the elegant paper by L.~Hatcher \cite{Hatcher}.
Answering the question posed in \cite{CMS} on the estimates of $\Phi(d,1,\O)$
in geometric terms, he showed that
$$
c_1(d) \,\frac{\mes_{d-1}(\dd\O)^d}{\mes_d (\O)^{d-1}}
\le \Phi(d,1,\O) \le 
c_2(d) \,\frac{\mes_{d-1}(\dd\O)^d}{\mes_d (\O)^{d-1}}
$$
for all convex domains $\O$.
The constants $c_1$ and $c_2$ depend only on $d$.
However, $\log c_1 (d)$ is of order of $d^2$ for large $d$
which is too large for us.

P.~Freitas and M.~Gama in \cite{FrG} showed that $\Psi (d,k,\O)$ 
grows as a power of $k$ as $k\to\infty$.
We are interested in another regime, 
namely what happens if $d \to \infty$.

Let us formulate our main result.

\begin{theorem}
\label{t14}
There is an absolute constant $C_{conv} > 0$ such that
\begin{equation*} 
\Psi (d,k,\O) \ge C_{conv} \left(\frac{e}2\right)^d 
\end{equation*}
for all natural $k$ and for all bounded convex domains $\O$.
\end{theorem}

Theorem \ref{t14} together with Corollary \ref{c12} invite us to make the following

\begin{conj}
\label{conj15}
There exist two absolute positive constants $C_*$, $\al_*$ such that
$$
\Psi (d,k,\O) \ge C_* e^{\al_* d} \qquad \forall \ k \in \N, \quad \forall \ \O \subset \R^d.
$$
\end{conj}

It is natural to consider the examples of multi-dimensional cubes and balls.
It is shown in \cite{CMS} that a) in the case of a ball the function $\Phi (d,1,B)$ 
grows faster than any power of $d$ (Theorem 3 in \cite{CMS});
and b) in the case of a cube it grows exponentially,
$\al_1 d \le \log \Phi (d,1,Q) \le \al_2 d$ 
with some positive constants $\al_1, \al_2$ (formula (10) in \cite{CMS}).
We refine these results by explicitly computing 
the asymptotics of $\log \Phi (d,1)$ in these cases.

\begin{theorem}
\label{t16}
Let $B$ be a $d$-dimensional ball.
Then
$$
\log \Phi (d, 1, B) = \al_{ball} d + O (d^{1/3}), \qquad d \to + \infty.
$$
Here
$$
\al_{ball} = \sqrt 2 \log (\sqrt 2 + 1) - \log 2 \approx 0,55.
$$
\end{theorem}

\begin{rem}
Moreover, one can show that
$$
\log \Phi (d, 1, B) = \al_{ball} d + \hat \al d^{1/3} + O (\log d),\qquad d \to + \infty,
$$
see \eqref{410} below.
\end{rem}

Note that this Theorem shows that if the Conjecture \ref{conj15} holds true
then $\al_* \le \al_{ball}$.

\begin{theorem}
\label{t17}
Let $Q$ be a $d$-dimensional cube.
Then the asymptotics
$$
\log \Phi (d, 1, Q) = \al_{cube} d + O (d^{1/3}), \qquad d \to + \infty
$$
holds true, $\al_{cube} \approx 1,04$.
\end{theorem}

We prove this Theorem in a separate paper.

For further references introduce the counting functions 
for the spectra of the Dirichlet and Neumann problems,
$$
N_{\cal D} (\O,\L) := \# \{ j: \la_j \le \L\}, \qquad
N_{\cal N} (\O,\L) := \# \{ j: \mu_j \le \L\} .
$$
In these terms
\begin{equation}
\label{14}
\Phi (d, k, \O) = N_{\cal N} (\O, \la_k).
\end{equation}
If $\la_k$ is a simple eigenvalue of the Dirichlet problem then
$$
\Psi (d, k, \O) = N_{\cal N} (\O, \la_k) - N_{\cal D} (\O, \la_k).
$$
In general case,
\begin{equation}
\label{15}
\Psi (d, k, \O) \ge N_{\cal N} (\O, \la_k) - N_{\cal D} (\O, \la_k).
\end{equation}

In the next section we prove Theorem \ref{t11} and Corollaries \ref{c12} and \ref{c13}.
In \S 3 we prove Theorem \ref{t14}.
In \S 4 we prove Theorem \ref{t16}.

{\bf Acknowledgements.} The results were partially obtained during the programme
{\it Geometric spectral theory and applications}, supported by EPSRC grant EP/Z000580/1.
The author thanks the Isaac Newton Institute for Mathematical Sciences, Cambridge, 
for support and hospitality.
The author also thanks Sergei V. Ivanov for consultation on the Kakeya-P\'al problem,
see \S \ref{smw}.

%%%%%%%%%%%%%%%%%%%%%%%%%%%%%%%%%%%%%%%%%%
\section{Bounds for $\Phi (d, 1, \O)$ and $\Phi (d, 2, \O)$ for arbitrary domains}

We use three following famous inequalities.

\begin{theorem}[Kr\" oger, \cite{Kr}]
\label{t21}
Let $\O \subset \R^d$ be a bounded domain such that
the embedding $W_2^1 (\O) \subset L_2 (\O)$ is compact.
Then the Neumann counting function satisfies the bound
$$
N_{\cal N} (\O,\la) \ge \frac2{d+2} \cdot \frac{\oo_d \mes_d \O}{(2\pi)^d}\ \la^{d/2}.
$$
\end{theorem}

Here and in what follows we denote by $\oo_d$ 
the volume of the unit ball in $\R^d$,
$$
\oo_d = \frac{\pi^{d/2}}{\G \left(\frac{d}2+1\right)}.
$$

\begin{theorem}[Faber-Krahn, \cite{Faber, Krahn}]
\label{t22}
Let $\O \subset \R^d$ be a bounded domain.
Then 
$$
\la_1 (\O) \ge \la_1 (B_R),
$$
where $B_R$ is a ball of the same volume, $\mes_d B_R = \mes_d \O$.
\end{theorem}

\begin{theorem}[Krahn-Szego, \cite{Kra}]
\label{t23}
Let $\O \subset \R^d$ be a bounded domain.
Then 
$$
\la_2 (\O) \ge \la_1 (B_{\tilde R}),
$$
where $B_{\tilde R}$ is a ball such that $\mes_d B_{\tilde R} = \frac12 \mes_d \O$.
\end{theorem}

{\it Proof of Theorem \ref{t11}.}
The first eigenfunction and the first eigenvalue 
of the Dirichlet problem in a ball of radius $R$ are well known:
\begin{equation}
\label{21}
\ph_1 (x) = |x|^{1-d/2} J_{\frac{d}2-1} \left(\frac{j_{\frac{d}2-1}|x|}R\right), 
\qquad
\la_1 (B_R) = \left(\frac{j_{\frac{d}2-1}}R\right)^2.
\end{equation}

The equalities $\mes_d B_R = \mes_d \O$ and $\mes_d B_{\tilde R} = \frac12 \mes_d \O$
imply
$$
R = \frac{(\mes_d \O)^{1/d}}{\oo_d^{1/d}}, \qquad
\tilde R = \frac{(\mes_d \O)^{1/d}}{(2 \oo_d)^{1/d}}.
$$
Therefore,
$$
\la_1 (\O) \ge \frac{j_{\frac{d}2-1}^2 \oo_d^{2/d}}{(\mes_d \O)^{2/d}},
\qquad 
\la_2 (\O) \ge \frac{j_{\frac{d}2-1}^2 2^{2/d} \oo_d^{2/d}}{(\mes_d \O)^{2/d}}.
$$
Now, Theorem \ref{t21} yields the bounds
\begin{eqnarray*}
\Phi (d, 1, \O) = N_{\cal N} (\O, \la_1) 
\ge \frac2{d+2} \cdot \frac{\oo_d \mes_d \O}{(2\pi)^d}\ \la_1 (\O)^{d/2} \\
\ge \frac{2 \oo_d^2 \left(j_{\frac{d}2-1}\right)^d}{(d+2) (2\pi)^d}
= \frac{\left(j_{\frac{d}2-1}\right)^d}{(d+2) 2^{d-1} \G(d/2+1)^2},
\end{eqnarray*}
and in the same manner,
\begin{eqnarray*}
\Phi (d, 2, \O) 
\ge \frac2{d+2} \cdot \frac{\oo_d \mes_d \O}{(2\pi)^d}\ \la_2 (\O)^{d/2}
\ge \frac{(j_{\frac{d}2-1})^d}{(d+2) 2^{d-2} \G(d/2+1)^2}.
\quad \qed
\end{eqnarray*}

{\it Proof of Corollary \ref{c12}.}
Taking a logarithm of inequality \eqref{12} we obtain 
\begin{equation}
\label{22}
\log \Phi (d,1,\O) \ge d \log j_{\frac{d}2-1} - d \log 2 - 2 \log \G\left(\frac{d}2 + 1\right)
+ O (\log d), \quad d \to \infty.
\end{equation}
The asymptotics of the zeros of the Bessel functions 
$$
j_\nu = \nu + |a_1| 2^{-1/3} \nu^{1/3} + O \left(\nu^{-1/3}\right), 
\qquad \nu \to + \infty,
$$
is well known, see for example \cite{QuWong}.
Here $a_1$ is the first zero of the Airy function $Ai$, $a_1 \approx - 2,34$.
Therefore,
\begin{equation}
\label{23}
j_{\frac{d}2-1} = \frac{d}2 + |a_1| 2^{-2/3} d^{1/3} + O(1), \quad d \to \infty,
\end{equation}
\begin{equation}
\label{24}
\log j_{\frac{d}2-1} = \log d - \log 2 + |a_1| 2^{1/3} d^{-2/3} + O(d^{-1}),
\quad d \to \infty.
\end{equation}
We use also Stirling's formula 
\begin{equation}
\label{25}
\log \G\left(\frac{d}2 + 1\right)
= \frac{d}2 \left(\log d - \log 2 - 1\right) + O(\log d), \quad d \to \infty.
\end{equation}
Substituting \eqref{24} and \eqref{25} into \eqref{22} we get
$$
\log \Phi (d,1,\O) \ge 
(1 - \log 2) d + |a_1| 2^{1/3} d^{1/3} + O(\log d)
\ge (1 - \log 2) d
$$
for sufficiently large $d$.
This yields the existence of such a constant $C_1 > 0$ that
$$
\Phi (d,1,\O) \ge C_1 \left(\frac{e}2\right)^d \qquad \qed
$$

For the proof of Corollary \ref{c13} we use

\begin{lemma}
\label{l21}
The function
$$
f_0 (x) = \log (1+x) - \frac{2x}3
$$
is positive on the interval $(0,1]$.
\end{lemma}

\begin{proof}
We have $f_0''(x) = \frac{-1}{(1+x)^2}$, thus, 
the function $f_0$ is concave.
Next, 
$$
f_0 (0) = 0, \qquad f_0(1) = \log 2 - \frac23 > 0 .
$$
Therefore,
$f_0(x) > 0$ if $0<x \le 1$.
\end{proof}

Denote the right hand side of \eqref{12} by
$$
f_1 (d) := \frac{(j_{d/2-1})^d}{(d+2) 2^{d-1} \G(d/2+1)^2}.
$$
From \cite{QuWong} we know that
$$
j_\nu \ge \nu + |a_1| 2^{-1/3} \nu^{1/3} ,
$$
and therefore
$$
j_{\nu-1} \ge (\nu - 1) \left(1  + |a_1| 2^{-1/3} \nu^{-2/3}\right).
$$ 
It is known also that
$$
\G (x+1) = x \G(x) < x^{x+1/2} e^{-x} \sqrt{2\pi} \,e^{\frac1{12x}},
$$
so,
$$
\G\left(\frac{d}2+1\right) < 
\sqrt{2\pi} \left(\frac{d}2\right)^{\frac{d+1}2} e^{-\frac{d}2 + \frac1{6d}}.
$$
Combining all together we obtain
$$
f_1(d) \ge \frac{(d-2)^d \left(1 + |a_1|2^{1/3} d^{-2/3}\right)^d e^d}
{\pi (d+2) 2^{d-1} d^{d+1} e^{\frac1{3d}}} =: f_2 (d).
$$

\begin{lemma}
\label{l22}
The function $f_2(d)$ is increasing for $d \ge 8$.
\end{lemma}

\begin{proof}
We have
\begin{eqnarray*}
\log f_2 (d) = d \log (d-2) + d \log \left(1 + |a_1|2^{1/3} d^{-2/3}\right) + d \\
- \log \pi - \log (d+2) - (d-1) \log 2 - (d+1) \log d - \frac1{3d},
\end{eqnarray*}
\begin{eqnarray*}
\left(\log f_2 (d)\right)' 
= \frac{d}{d-2} + \log (d-2) + \log \left(1 + |a_1|2^{1/3} d^{-2/3}\right)
- \frac{|a_1| 2^{4/3}}{3 \left(1 + |a_1|2^{1/3} d^{-2/3}\right) d^{2/3}} \\
+1 - \frac1{d+2} - \log 2 - \frac{d+1}d - \log d + \frac1{3d^2}.
\end{eqnarray*}
Furthermore,
$$
\frac{d}{d-2} - \frac1{d+2} - \frac{d+1}d = \frac{6d+4}{d(d^2-4)} > 0,
$$
$$
\log \left(1 + |a_1|2^{1/3} d^{-2/3}\right)
- \frac{|a_1| 2^{4/3}}{3 \left(d^{2/3} + |a_1|2^{1/3}\right)}
> \log \left(1 + |a_1|2^{1/3} d^{-2/3}\right)
- \frac{|a_1| 2^{4/3}}{3 d^{2/3}} > 0
$$ 
by virtue of Lemma \ref{l21} as we have $|a_1| 2^{1/3} d^{-2/3} \le 1$
if $d \ge \sqrt 2 |a_1|^{3/2} \approx 5,06$.
Thus,
$$
\left(\log f_2 (d)\right)' \ge
\log (d-2) + 1 - \log 2 - \log d = \log \left(\frac{(d-2) e}{2d}\right) > 0,
$$
whenever
$d > \frac{2e}{e-2} \approx 7,57$.
\end{proof}

{\it Proof of Corollary \ref{c13}.}
The preceding Lemma yields
\begin{equation}
\label{26}
f_1 (d) \ge f_2 (d) \ge f_2 (11) \approx 2,37 \qquad
\text{if} \quad d \ge 11.
\end{equation}
For any $\er \in (0,1)$ denote by $\tilde \er (\O)$ the number
$$
\tilde \er (\O) =  \frac{j_{\frac{d}2-1}^2 \oo_d^{2/d}}{(\mes_d \O)^{2/d}}
\left(1 - (1-\er)^{2/d}\right) > 0.
$$
Then repeating the arguments from the proof of Theorem \ref{t11} for 
the counting function at the point $(\la_1 (\O) - \tilde \er(\O))$ one obtain
$$
N_{\cal N} (\O, \la_1 - \tilde \er (\O)) 
\ge \frac2{d+2} \cdot \frac{\oo_d \mes_d \O}{(2\pi)^d} 
\left(\frac{j_{\frac{d}2-1}^2 \oo_d^{2/d}}{(\mes_d \O)^{2/d}} - \tilde \er (\O)\right)^{\frac{d}2}
= f_1 (d) (1 - \er).
$$
For sufficiently small $\er$ the inequality \eqref{26} implies
$$
N_{\cal N} (\O, \la_1 - \tilde \er) \ge 2,3 \quad \Rightarrow
\quad \mu_3 \le \la_1 - \tilde \er < \la_1
$$
whenever $d \ge 11$.
In the same manner,
$$
\mu_5 < \la_2 \qquad \text{whenever} \quad d \ge 11.
$$
For small values of $d$ we calculate the function $f_1(d)$ explicitly:
$$
\begin{array}{|c|c|c|c|c|}
\hline
& & & &\\
d & j_{\frac{d}2-1} & f_1(d) & \text{inequality with}\ \la_1 &
\text{inequality with}\ \la_2 \\
\hline
& & & &\\
6 & 5,1356 & 1,99 & & \mu_4 < \la_2 \\
\hline
& & & &\\
7 & 5,7635 & 2,71 & \mu_3 < \la_1 & \mu_6 < \la_2\\
\hline
& & & &\\
8 & 6,3802 & 3,72 & \mu_4 < \la_1 & \mu_8 < \la_2\\
\hline
& & & &\\
9 & 6,9879 & 5,15 & \mu_6 < \la_1 & \mu_{11} < \la_2\\
\hline
& & & &\\
10 & 7,5883 & 7,16 & \mu_8 < \la_1 & \mu_{15} < \la_2\\
\hline
\end{array}
$$
\qed

%%%%%%%%%%%%%%%%%%%%%%%%%%%%%%%%%%%%%%%%%%
\section{Convex domains}
In this section we prove Theorem \ref{t14}.

\subsection{Estimate for the difference $N_{\cal N} (\la) - N_{\cal D} (\la)$
via the exponentials}
The following Lemma is proved in \cite{FilSaf}.
We give a proof for the sake of completeness.

\begin{lemma}
\label{l31}
Let $\la > 0$.
Introduce the subspace
\begin{equation}
\label{30}
H_{\cal D} (\la) = \operatorname{Lin} \left\{\ph_1, \dots, \ph_k\right\}_{\la_k \le \la},
\end{equation}
where $\ph_j$ is the Dirichlet eigenfunction corresponding to $\la_j$.
Let $F \subset W_2^1 (\O)$ be a subspace such that
$$
F \cap H_{\cal D} = \{0\}, \qquad
- \D v = \la v, \qquad
\|\n v\|_{L_2(\O)}^2 \le \la \|v\|_{L_2(\O)}^2, \quad \forall \ v \in F.
$$
Then
$$
N_{\cal N} (\O, \la) - N_{\cal D} (\O, \la) \ge \dim F.
$$
\end{lemma}

\begin{proof}
Let us consider the subspace $H_{\cal D} + F \subset W_2^1 (\O)$.
By assumptions,
$$
\dim \left(H_{\cal D} + F\right) = N_{\cal D} (\O, \la) + \dim F.
$$
Pick a function $w$ from this subspace,
$$
w = u + v, \quad u \in H_{\cal D}, \quad v \in F.
$$
Then we have
\begin{eqnarray*}
\int_\O |\n w|^2 dx = 
\int_\O \left(|\nabla u|^2 + |\nabla v|^2 +
2\, \re \langle \nabla u, \nabla v\rangle \right) dx, \\
\int_\O |\n u|^2 dx \le \la \int_\O |u|^2 dx, \quad
\int_\O |\n v|^2 dx \le \la \int_\O |v|^2 dx, \\
\int_\O \langle \nabla u, \nabla v\rangle 
= - \int_\O u \D \overline v \, dx = \la \int_\O u \overline v \, dx.
\end{eqnarray*}
Therefore,
$$
\int_\O |\n w|^2 dx
\le \la \int_\O\left(|u|^2 + |v|^2 + 2 \re (u \overline v)\right) dx
= \la \int_\O |w|^2 dx.
$$
Now, the variational principle implies
$$
N_{\cal N} (\O, \la) \ge \dim \left(H_{\cal D} + F\right).
\qquad \qedhere
$$
\end{proof}

Taking into account this Lemma and the inequality \eqref{15},
our goal now is to construct a such subspace $F$ for $\la \ge \la_1 (\O)$.

\subsection{Minimal width}
\label{smw}
Let $\O \subset \R^d$ be a bounded convex domain.
The distance between two supporting hyperplanes of $\O$
orthogonal to a vector $\vec e \in \R^d$ is called the width
of $\O$ in the direction $\vec e$.
We denote by $w_{\min} (\O)$ the smallest width of $\O$ 
over all directions.

\begin{lemma}
\label{l32}
Let $\O \subset \R^d$ be a bounded convex domain.
For each $k = 1, \dots, d$, there exists a $k$-dimensional plane $H_k$,
$\dim H_k = k$, such that
\begin{equation}
\label{31}
\mes_k (P_k) \ge \frac{w_{\min} (\O)^k}{k!},
\qquad \text{where}\quad P_k = \O \cap H_k.
\end{equation}
\end{lemma}

\begin{proof}
We prove this Lemma by induction in $k$.

Clearly, one can find two points $x_1^{(1)}, x_2^{(1)} \in \overline \O$
such that $|x_1^{(1)} - x_2^{(1)}| \ge w_{\min} (\O)$.
Now, take the line passing through these two points as $H_1$.
Thus, the claim in the case $k=1$ is proven.

Assume that $H_k$ satisfies \eqref{31}.
Choose a unit vector $\vec e_{k+1}$ orthogonal to the plane $H_k$.
We take the linear hull of $H_k$ and $\vec e_{k+1}$ as the next plane $H_{k+1}$,
$$
H_{k+1} = \left\{x = y + t \vec e_{k+1},\ y \in H_k, t \in \R\right\}.
$$
By definition of the minimal width of $\O$,
there are two points $x_1^{(k+1)}, x_2^{(k+1)} \in \overline \O$ such that
$$
x_1^{(k+1)} = y_1 + t_1 \vec e_{k+1}, \quad
x_2^{(k+1)} = y_2 - t_2 \vec e_{k+1}, \quad
y_1, y_2 \in H_k, \quad t_1, t_2 \ge 0, \quad t_1 + t_2 \ge w_{\min} (\O).
$$
Then the convex set $\overline{P_{k+1}} = \overline \O \cap H_{k+1}$
contains the convex hull of the set $P_k$ and of both points $x_1^{(k+1)}$, $x_2^{(k+1)}$.
The measure of this convex hull is at least
$$
\frac{t_1 \mes_k P_k}{k+1} + \frac{t_2 \mes_k P_k}{k+1}. 
$$
Now, due to \eqref{31}
$$
\mes_{k+1} P_{k+1} \ge \frac{(t_1+ t_2) \mes_k P_k}{k+1} 
\ge \frac{w_{\min} (\O)^{k+1}}{(k+1)!},
$$
which is \eqref{31} for the index $k+1$.
\end{proof}

Choosing $k=d$ in \eqref{31} we obtain

\begin{theorem}
\label{t33}
Let $\O \subset \R^d$ be a bounded convex domain.
Put
\begin{equation}
\label{32}
w \equiv w(\O) := \left(d! \mes_d \O\right)^{1/d}.
\end{equation}
Then
\begin{equation}
\label{33}
w_{\min} (\O) \le w (\O).
\end{equation}
\end{theorem}

\begin{rem}
We provided the proof of Theorem \ref{t33} for the sake of completeness.
The question about minimizing the volume of convex domain of given minimal width in $\R^d$
is known as Kakeya-P\'al problem.
The answer is known \cite{Pal} for the case $d=2$ only:
the optimal domain is a regular triangle of height $w$.
This fact yields the following improvement of the estimate \eqref{33}:
$$
w_{\min} (\O)^d \le \frac{\sqrt 3}2\, d! \mes_d \O.
$$
It was done in \cite{Firey}.
Another improvement is contained in \cite[the proof of Theorem 6.2]{Bezdek}.
\end{rem}

In Theorem \ref{t14} we assume that the domain $\O$ is convex.
The only advantage of it we use is the fact that $\O$ is contained in a layer
of width $w(\O)$.

\begin{theorem}
\label{t335}
Assume that a bounded domain $\O \subset \R^d$ is contained in a layer
of width $w(\O)$ defined in \eqref{32}.
Then 
\begin{equation*} 
\Psi (d,k,\O) \ge C_{conv} \left(\frac{e}2\right)^d 
\end{equation*}
for all $k$.
Here $C_{conv}>0$ is an absolute constant.
\end{theorem}

Clearly, Theorem \ref{t14} follows from Theorem \ref{t33} and Theorem \ref{t335}.

\subsection{Functions $u_\eta$}
In the sequel in this section we assume
$$
x = (x_1, y) \in \R^d, \qquad x_1 \in \R, \quad y \in \R^{d-1},
$$
$$
\O \subset\left\{x = (x_1, y) : |x_1| \le \frac{w}2\right\},
$$
where the number $w = w(\O)$ is defined by the formula \eqref{32}.

\begin{lemma}
\label{l36}
Let $\de > 0$.
Then
$$
\int_\O \cos (2 x_1 \de) dx \ge \left(1 - \frac{w^2\de^2}2\right) \mes_d \O.
$$
\end{lemma}

\begin{proof}
The inequality
$$
0 \le 1 - \cos z \le \frac{z^2}2 \qquad \forall \ z \in \R
$$
implies the bound
$$
\left| \mes_d \O - \int_\O \cos (2 x_1 \de) dx\right|
\le \int_\O \left|1-\cos (2 x_1 \de)\right| dx
\le \int_\O 2 x_1^2 \de^2 dx
\le \frac{w^2\de^2}2\, \mes_d \O.
\quad \qedhere
$$
\end{proof}

Let
$$
\la \ge \de^2, \qquad \eta \in \R^{d-1}, \qquad |\eta|^2 + \de^2 = \la.
$$
Introduce the function
\begin{equation}
\label{34}
u_\eta (x) = \cos (x_1 \de) \,e^{i\<y,\eta\>},
\qquad \text{where} \quad x = (x_1, y).
\end{equation}
It is clear that
\begin{equation}
\label{35}
- \D u_\eta = \la u_\eta.
\end{equation}

\begin{lemma}
\label{l37}
Let
$$
\eta, \ze \in \R^{d-1}, \qquad
|\eta| = |\ze| = \sqrt{\la - \de^2}.
$$
Then
$$
\la \int_\O u_\eta\, \overline{u_\ze}\, dx - \int_\O \<\n u_\eta, \n u_\ze\> dx
= \int_\O \left(\de^2 \cos(2x_1\de) + \frac{|\eta-\ze|^2}2 \cos^2 (x_1\de)\right)
e^{i\<y, \eta-\ze\>} dx.
$$
\end{lemma}

\begin{proof}
We have
$$
\int_\O u_\eta\, \overline{u_\ze}\, dx
= \int_\O \cos^2 (x_1\de) e^{i\<y, \eta-\ze\>} dx,
$$
$$
\int_\O \<\n u_\eta, \n u_\ze\> dx
= \int_\O \left(\de^2 \sin^2 (x_1 \de) + \<\eta, \ze\> \cos^2 (x_1 \de)\right) e^{i\<y, \eta-\ze\>} dx.
$$
Taking into account the identity
$$
\frac{|\eta-\ze|^2}2 = \la - \de^2 - \<\eta,\ze\>,
$$
we obtain the result.
\end{proof}

Lemma \ref{l36} and Lemma \ref{l37} imply the following

\begin{cor}
\label{c38}
Let $\eta \in \R^{d-1}$, $|\eta| = \sqrt{\la - \de^2}$.
Then
$$
\la \int_\O |u_\eta|^2 dx - \int_\O |\n u_\eta|^2 dx 
\ge \de^2 \left(1 - \frac{w^2\de^2}2\right) \mes_d \O.
$$
\end{cor}

In the sequel we take
$$
\de = \frac1{w(\O)}.
$$
Then

\begin{equation}
\label{36}
\la \int_\O |u_\eta|^2 dx - \int_\O |\n u_\eta|^2 dx 
\ge \frac{\mes_d \O}{2 w(\O)^2}.
\end{equation}

\begin{rem}
In the sequel we consider 
$$
\la \ge \la_1 (\O) \ge \frac{j_{\frac{d}2-1}^2 \oo_d^{2/d}}{(\mes_d \O)^{2/d}},
\qquad \text{and} \qquad \de^2 = \frac1{\left(d! \mes_d\O\right)^{2/d}}.
$$
Note that for each $d \ge 2$ we have
$j_{\frac{d}2-1} \ge j_0 > 1$, and
$$
\oo_d = \frac{\pi^{d/2}}{\G \left(\frac{d}2+1\right)} 
> \frac1{\G \left(\frac{d}2+1\right)} > \frac1{d!},
$$ 
so the condition $\la \ge \de^2$ is fulfilled.
\end{rem}

\subsection{Matrix $B(H)$}
We consider the set of $N$ points on the sphere $S_{\sqrt{\la-\de^2}}^{d-2}$,
$$
H := \left\{\eta_1, \dots, \eta_N\right\} \in \left(S_{\sqrt{\la-\de^2}}^{d-2}\right)^N.
$$
The number $N$ will be chosen later.
For each such set $H$ we define the linear space
$$
F (H) = \operatorname{Lin} \left\{u_{\eta_1}, \dots, u_{\eta_N}\right\},
$$
the functions $u_\eta$ here are defined by the formula \eqref{34}.
Clearly,
$$
F (H) \subset W_2^1 (\O).
$$
If all points $\eta_1, \dots, \eta_N$ are different then $\dim F(H) = N$.
Denote by $\Sigma^{(1)} \subset \left(S_{\sqrt{\la-\de^2}}^{d-2}\right)^N$
the set of such sets $H$ for which
\begin{equation}
\label{375}
F (H) \cap H_{\cal D} (\la) = \{0\},
\end{equation}
the space $H_{\cal D} (\la)$ is defined by \eqref{30}.

\begin{lemma}
\label{l39}
The set $\Sigma^{(1)}$ is everywhere dense in $\left(S_{\sqrt{\la-\de^2}}^{d-2}\right)^N$.
\end{lemma}

\begin{proof}
Pick an arbitrary set
$$
\left\{\eta_1, \dots, \eta_N\right\} \in \left(S_{\sqrt{\la-\de^2}}^{d-2}\right)^N.
$$
All functions $\{u_\eta\}_{\eta \in S_{\sqrt{\la-\de^2}}^{d-2}}$ are linearly independent.
Therefore in an arbitrarily small neighborhood of the point $\eta_1$
one can find a point $\tilde \eta_1$ such that
$u_{\tilde \eta_1} \notin H_{\cal D} (\la)$.
In any neighborhood of the point $\eta_2$ one can find a point $\tilde \eta_2$ such that
$$
u_{\tilde \eta_2} \notin \operatorname{Lin} \left\{u_{\tilde \eta_1}, \ph_1, \dots, \ph_k\right\}.
$$
In any neighborhood of the point $\eta_3$ one can find a point $\tilde \eta_3$ such that
$$
u_{\tilde \eta_3} \notin 
\operatorname{Lin} \left\{u_{\tilde \eta_1}, u_{\tilde \eta_2}, \ph_1, \dots, \ph_k\right\},
$$
and so on.
For the final set $\{\tilde\eta_1, \dots, \tilde\eta_N\}$ all the functions
$u_{\tilde \eta_1}, \dots, u_{\tilde \eta_N}, \ph_1, \dots, \ph_k$
are linearly independent, and therefore,
$$
F \left(\left\{\tilde\eta_1, \dots, \tilde\eta_N\right\}\right) \cap H_{\cal D} (\la) = \{0\},
$$
and 
$\left\{\tilde\eta_1, \dots, \tilde\eta_N\right\} \in \Sigma^{(1)}$.
\end{proof}

If 
$$
f = \sum_{j=1}^N c_j u_{\eta_j} \in F (H), \qquad c_j \in \C,
$$
then by \eqref{35}
$$
- \D f(x) = \la f(x),
$$
and
$$
\la \int_\O |f(x)|^2 dx - \int_\O |\n f(x)|^2 dx = (B \vec c, \vec c).
$$
Due to Lemma \ref{l37} the entries of the $(N\times N)$-matrix $B$ equal to
\begin{eqnarray*}
B_{kj} =
\la \int_\O u_{\eta_j}\, \overline{u_{\eta_k}}\, dx 
- \int_\O \<\n u_{\eta_j}, \n u_{\eta_k}\> dx \\
= \int_\O \left(\de^2 \cos(2x_1\de) + \frac{|\eta_j-\eta_k|^2}2 \cos^2 (x_1\de)\right)
e^{i\<y, \eta_j-\eta_k\>} dx.
\end{eqnarray*}
Represent the matrix $B$ as a sum of diagonal and off-diagonal parts:
$$
B = B_0 + B_1, \qquad B_0 = \operatorname{diag} (B_{jj}).
$$
By virtue of Corollary \ref{c38} and formula \eqref{36} we have
\begin{equation}
\label{38}
B_0 \ge \frac{\mes_d \O}{2 w(\O)^2}\, I_{N\times N}.
\end{equation}
Let us estimate the Hilbert-Schmidt norm of the matrix $B_1$.
We have
\begin{eqnarray*}
|B_{kj}|^2 
= \left|\de^2 \int_\O \cos(2x_1\de) e^{i\<y, \eta_j-\eta_k\>} dx
+ \frac{|\eta_j-\eta_k|^2}2 \int_\O \cos^2 (x_1\de) e^{i\<y, \eta_j-\eta_k\>} dx\right|^2 \\
\le \left(\de^4 + \frac{|\eta_j-\eta_k|^4}4\right)
\left(G_1 (\eta_j-\eta_k) + G_2 (\eta_j-\eta_k)\right),
\end{eqnarray*}
where we used the notation
\begin{equation}
\label{39}
G_1 (\te) = \left|\int_\O \cos(2x_1\de) e^{i\<y,\te\>} dx\right|^2, \quad
G_2 (\te) = \left|\int_\O \cos^2 (x_1\de) e^{i\<y,\te\>} dx\right|^2,
\end{equation}
and the inequality $|ab+cd|^2 \le (|a|^2+|c|^2) (|b|^2 + |d|^2)$.
Therefore,
$$
\|B_1\|_{S_2}^2 = \sum_{j\neq k} |B_{kj}|^2
\le \sum_{j \neq k} \left(\de^4 + \frac{|\eta_j-\eta_k|^4}4\right)
\left(G_1 (\eta_j-\eta_k) + G_2 (\eta_j-\eta_k)\right).
$$
Now, integrate this inequality over all $H \in \left(S_{\sqrt{\la-\de^2}}^{d-2}\right)^N$:
\begin{eqnarray}
\nonumber
\int_{S_{\sqrt{\la-\de^2}}^{d-2}} \dots \int_{S_{\sqrt{\la-\de^2}}^{d-2}} 
\|B_1 (H)\|_{S_2}^2 dS (\eta_1) \dots dS (\eta_N) \\
\label{310}
\le N (N-1) \left(\mes_{d-2} S_{\sqrt{\la-\de^2}}^{d-2}\right)^{N-2} \times \\
\times \int_{S_{\sqrt{\la-\de^2}}^{d-2}} \int_{S_{\sqrt{\la-\de^2}}^{d-2}}
\left(\de^4 + \frac{|\eta-\ze|^4}4\right)
\left(G_1 (\eta-\ze) + G_2 (\eta-\ze)\right) dS (\eta) dS(\ze).
\nonumber
\end{eqnarray}

\subsection{Estimates of integrals}
The following identity is proved in \cite[Lemma 5.2]{FilSaf}.

\begin{lemma}
Let $m \ge 2$, $R>0$, $f \in C(\R^n)$.
Then
\begin{eqnarray*}
\int_{S_R^{m-1}} \int_{S_R^{m-1}} f (\eta-\ze) dS (\eta) dS (\ze) 
= (m-1)\, \oo_{m-1} R^2 \int_{B_{2R}^m} f(\te) \left(R^2 - \frac{\te^2}4\right)^{\frac{m-3}2}
\frac{d\te}{|\te|} .
\end{eqnarray*}
\end{lemma}

This Lemma yields the equality
\begin{eqnarray}
\label{311}
\int_{S_{\sqrt{\la-\de^2}}^{d-2}} \int_{S_{\sqrt{\la-\de^2}}^{d-2}}
\left(\de^4 + \frac{|\eta-\ze|^4}4\right)
\left(G_1 (\eta-\ze) + G_2 (\eta-\ze)\right) dS (\eta) dS(\ze) \\
= (d-2) \oo_{d-2} (\la - \de^2) 
\int_{B_{2\sqrt{\la-\de^2}}^{d-1}} \left(\de^4 + \frac{|\te|^4}4\right)
\left(G_1(\te) + G_2(\te)\right) \left(\la - \de^2 - \frac{\te^2}4\right)^{\frac{d-4}2} \frac{d\te}{|\te|} .
\nonumber
\end{eqnarray}

\begin{lemma}
\label{l311}
The estimates
$$
\int_{\R^{d-1}} \left|\int_\O \cos(2x_1\de) e^{i\<y,\te\>} dx\right|^2 d\te
\le (2\pi)^{d-1} w(\O) \mes_d \O,
$$
$$
\int_{\R^{d-1}} \left|\int_\O \cos^2 (x_1\de) e^{i\<y,\te\>} dx\right|^2 d\te
\le (2\pi)^{d-1} w(\O) \mes_d \O
$$
hold true.
\end{lemma}

\begin{proof}
Given $x_1$ denote by $\O' (x_1)$ a cross-section of the domain $\O$ 
with the corresponding hyperplane,
$$
\O' (x_1) = \left\{y \in \R^{d-1} : (x_1, y) \in \O\right\}.
$$
Then
$$
\left\|\int_{\O'(x_1)} e^{i\<y,\te\>} dy\right\|_{L_2(\R^{d-1})}^2 
= \int_{\R^{d-1}} \left|\int_{\O'(x_1)} e^{i\<y,\te\>} dy\right|^2 d\te
= (2\pi)^{d-1} \mes_{d-1} \O' (x_1),
$$
as the Fourier transform preserves $L_2$-norm.
Therefore,
\begin{eqnarray*}
\left\|\int_\O \cos (2x_1\de) e^{i\<y,\te\>} dx\right\|_{L_2(\R^{d-1})}
= \left\|\int_{-w/2}^{w/2} \cos (2x_1\de) dx_1 
\int_{\O'(x_1)} e^{i\<y,\te\>} dy\right\|_{L_2(\R^{d-1})} \\
\le \int_{-w/2}^{w/2} dx_1 \left\|\int_{\O'(x_1)} e^{i\<y,\te\>} dy\right\|_{L_2(\R^{d-1})} 
= \int_{-w/2}^{w/2} dx_1 (2\pi)^{\frac{d-1}2} \left(\mes_{d-1} \O' (x_1)\right)^{1/2} \\
\le (2\pi)^{\frac{d-1}2} \sqrt{w} \left(\int_{-w/2}^{w/2} dx_1 \mes_{d-1} \O' (x_1)\right)^{1/2}
= (2\pi)^{\frac{d-1}2} \sqrt{w} \left(\mes_d \O\right)^{1/2}.
\end{eqnarray*}
The proof of the second inequality is similar.
\end{proof}

\begin{lemma}
\label{l312}
Let $m \ge 2$, $a,b \ge 0$, $G \in C(\R^m)$,
$$
0 \le G(\te) \le a \quad \forall \ \te \in \R^m, 
\qquad \int_{\R^m} G(\te) \, d\te \le b.
$$
Then
$$
\int_{\R^m} \frac{G(\te)\, d\te}{|\te|} \le 
\frac{m}{m-1} \left(\oo_m a b^{m-1}\right)^{1/m}.
$$
\end{lemma}

\begin{proof}
Let us consider the function
$$
h(r) = \int_{|\te| \ge r} \frac{G(\te)\, d\te}{|\te|}.
$$
Clearly, the integral $\int_{\R^m} \frac{G(\te)\, d\te}{|\te|}$ converges,
thus, the function $h$ is well defined, non-negative and continuous on $[0,\infty)$.
Moreover,
$$
r h(r) \le \int_{|\te| \ge r} G(\te)\, d\te, \quad \text{so} \quad
r h(r) \mathop{\longrightarrow}\limits_{r \to \infty} 0.
$$
Next,
\begin{equation}
\label{312}
h'(r) = - \frac1r \int_{|\te|=r} G(\te) \, dS(\te),
\end{equation}
therefore,
$$
\int_{\R^m} G(\te) \, d\te = \int_0^\infty dr \int_{|\te|=r} G(\te) \, dS(\te)
= - \int_0^\infty r h'(r)\, dr = \int_0^\infty h(r) \, dr,
$$
and
\begin{equation}
\label{313}
\int_0^\infty h(r) \, dr \le b.
\end{equation}
Furthermore, \eqref{312} implies
$$
h'(r) \ge - m \oo_m a r^{m-2},
$$
so,
$$
h(r) = h(0) + \int_0^r h'(t)\,dt
\ge h(0) - m \oo_m a \int_0^r t^{m-2} dt 
= h(0) - \frac{m \oo_m a}{m-1} \, r^{m-1}.
$$
Denote
$$
r_0 = \left(\frac{(m-1) h(0)}{m \oo_m a}\right)^{\frac1{m-1}}.
$$
Then by virtue of \eqref{313}
\begin{eqnarray*}
b \ge \int_0^\infty h(r)\, dr 
\ge \int_0^{r_0} \left(h(0) - \frac{m \oo_m a}{m-1} \, r^{m-1}\right) dr \\
= h(0) r_0 - \frac{\oo_m a r_0^m}{m-1}
= \frac{m-1}m \, r_0 h(0)
= \left(\frac{(m-1) h(0)}m\right)^{\frac{m}{m-1}} \frac1{(\oo_m a)^{\frac1{m-1}}},
\end{eqnarray*}
and therefore,
$$
\left(\oo_m a b^{m-1}\right)^{1/m} \ge \frac{m-1}m \, h(0).
\qquad \qedhere
$$
\end{proof}

\begin{rem}
The constant in this estimate is sharp.
Indeed, if we consider the function
$$
G(\te) = \begin{cases}
1, \quad &|\te| < 1, \\
0, \quad &|\te| \ge 1,
\end{cases}
$$
then
$$
a = 1, \quad b = \oo_m \quad \text{and} \quad
\int_{\R^m} \frac{G(\te)\, d\te}{|\te|} = \frac{m \oo_m}{m-1}.
$$
\end{rem}

\begin{cor}
\label{c314}
Functions $G_1$, $G_2$ defined in \eqref{39} satisfy the estimates
$$
\int_{\R^{d-1}} \frac{G_j(\te)\, d\te}{|\te|} \le 
\frac{(d-1) (2\pi)^{d-2}}{d-2} 
\left(\oo_{d-1} w(\O)^{d-2} \left(\mes_d\O\right)^d\right)^{\frac1{d-1}},
\quad j = 1, 2.
$$
\end{cor}

\begin{proof}
It is clear by definition that
$$
0 \le G_j (\te) \le (\mes_d\O)^2, \quad j = 1, 2.
$$
So, the functions $G_1$, $G_2$ satisfy the assumptions of Lemma \ref{l312}
with 
$$
m=d-1, \quad a = (\mes_d\O)^2 \quad \text{and}
\quad b =  (2\pi)^{d-1} w(\O) \mes_d \O
$$ 
due to Lemma \ref{l311}.
\end{proof}

\begin{lemma}
\label{l315}
Let $d\ge 5$.
Then
$$
x^2 (1-x)^{\frac{d-4}2} \le \frac{16 (d-4)^{\frac{d-4}2}}{d^{\frac{d}2}},
\qquad x \in [0,1].
$$
\end{lemma}

\begin{proof}
The function
$g(x): = x^2 (1-x)^{\frac{d-4}2}$
vanishes at the endpoints of the interval,
$g(0) = g(1) = 0$,
and it is positive inside.
The derivative
$$
g'(x) = 2x (1-x)^{\frac{d-4}2} - \frac{d-4}2\, x^2 (1-x)^{\frac{d-6}2}
$$
has a single root $x = \frac4d$.
Therefore,
$$
g(x) \le g\left(\frac4d\right) = \frac{16 (d-4)^{\frac{d-4}2}}{d^{\frac{d}2}}.
\qquad \qedhere
$$
\end{proof}

\begin{cor}
\label{c316}
If $\de^2 + \frac{\te^2}4 \le \la$ then
$$
\left(\de^4 + \frac{|\te|^4}4\right) \left(\la - \de^2 - \frac{\te^2}4\right)^{\frac{d-4}2}
\le \frac{16 (d-4)^{\frac{d-4}2}}{d^{\frac{d}2}}\, \la^{\frac{d}2}.
$$
\end{cor}

\begin{proof}
Clearly,
$$
\left(\de^4 + \frac{|\te|^4}4\right) \left(\la - \de^2 - \frac{\te^2}4\right)^{\frac{d-4}2}
\le \left(\de^2 + \frac{\te^2}4\right)^2 \left(\la - \de^2 - \frac{\te^2}4\right)^{\frac{d-4}2}.
$$
Applying the preceding Lemma with 
$x = \frac{\de^2 + \te^2/4}\la$ to the right hand side we get the claim.
\end{proof}

Now, we are ready to estimate the integral in the right hand side of \eqref{311}:
\begin{eqnarray}
\nonumber
\int_{B_{2\sqrt{\la-\de^2}}^{d-1}} \left(\de^4 + \frac{|\te|^4}4\right)
\left(G_1(\te) + G_2(\te)\right) \left(\la - \de^2 - \frac{\te^2}4\right)^{\frac{d-4}2} \frac{d\te}{|\te|} \\
\label{314}
\le \frac{16 (d-4)^{\frac{d-4}2}}{d^{\frac{d}2}}\, \la^{\frac{d}2}
\int_{\R^{d-1}} \left(G_1(\te) + G_2(\te)\right) \frac{d\te}{|\te|} \\
\le \frac{32 (d-4)^{\frac{d-4}2} \la^{\frac{d}2} (d-1) (2\pi)^{d-2}}{d^{\frac{d}2} (d-2)}
\left(\oo_{d-1} w(\O)^{d-2} \left(\mes_d\O\right)^d\right)^{\frac1{d-1}}.
\nonumber
\end{eqnarray}
Here we used Corollary \ref{c316} on the first step and Corollary \ref{c314} on the second step.
Relations \eqref{310}, \eqref{311} and \eqref{314} yield 
\begin{eqnarray}
\nonumber
\int_{S_{\sqrt{\la-\de^2}}^{d-2}} \dots \int_{S_{\sqrt{\la-\de^2}}^{d-2}} 
\|B_1 (H)\|_{S_2}^2 dS (\eta_1) \dots dS (\eta_N) 
\le N (N-1) \left(\mes_{d-2} S_{\sqrt{\la-\de^2}}^{d-2}\right)^{N-2} \times \\
\times 32 (d-4)^{\frac{d-4}2} (d-1) d^{-\frac{d}2} \oo_{d-2} (2\pi)^{d-2} (\la-\de^2) \la^{\frac{d}2}
\left(\oo_{d-1} w(\O)^{d-2} \left(\mes_d\O\right)^d\right)^{\frac1{d-1}}.
\label{315}
\end{eqnarray}
The function $\|B_1 (H)\|_{S_2}^2$ is a continuous function of 
$H \in \left(S_{\sqrt{\la-\de^2}}^{d-2}\right)^N$.
Therefore, there is a non-empty open subset 
$$
\Sigma^{(2)} \subset \left(S_{\sqrt{\la-\de^2}}^{d-2}\right)^N
$$
where this function does not exceed the right hand side of \eqref{315}
divided by the volume $\left(\mes_{d-2} S_{\sqrt{\la-\de^2}}^{d-2}\right)^N$.
Taking into account that
$$
\mes_{d-2} S_{\sqrt{\la-\de^2}}^{d-2}
= (d-1) \oo_{d-1} (\la-\de^2)^{\frac{d-2}2},
$$
we obtain for $H \in \Sigma^{(2)}$
\begin{eqnarray}
\nonumber
\|B_1 (H)\|_{S_2}^2 
\le \frac{N (N-1)}{\left(\mes_{d-2} S_{\sqrt{\la-\de^2}}^{d-2}\right)^2} \times \\
\label{317}
\times 32 (d-4)^{\frac{d-4}2} (d-1) d^{-\frac{d}2} \oo_{d-2} (2\pi)^{d-2} (\la-\de^2) \la^{\frac{d}2}
\left(\oo_{d-1} w(\O)^{d-2} \left(\mes_d\O\right)^d\right)^{\frac1{d-1}} \\
= \frac{32 N (N-1) (d-4)^{\frac{d-4}2} \oo_{d-2} (2\pi)^{d-2} \la^{\frac{d}2}}
{(d-1) d^{\frac{d}2} \oo_{d-1}^2 (\la-\de^2)^{d-3}}  
\left(\oo_{d-1} w(\O)^{d-2} \left(\mes_d\O\right)^d\right)^{\frac1{d-1}} .
\nonumber
\end{eqnarray}
The set $\Sigma^{(2)}$ is open, and by Lemma \ref{l39} 
the set $\Sigma^{(1)}$ is everywhere dense.
Therefore, 
$$
\Sigma^{(1)} \cap \Sigma^{(2)} \neq \emptyset.
$$

We proved the following

\begin{lemma}
\label{l317}
For any natural $N$ there is $H = H^{(N)} \in \left(S_{\sqrt{\la-\de^2}}^{d-2}\right)^N$
such that the relations
\eqref{375} and \eqref{317} are fulfilled.
\end{lemma}

\subsection{Positive spectra of matrices of special kind}
Given a Hermitian matrix $B \in \mat (\C, N \times N)$
denote by $N_+ (B)$ the number of its positive eigenvalues,
$N_+(B) = \# \left(\spec (B) \cap (0, \infty)\right)$.
The number $N_+ (B)$ coincides with the maximal dimension
of a subspace of $\C^N$ where the quadratic form of the matrix $B$ is positive.

\begin{lemma}
\label{l318}
Let 
$$
B = B^* \in \mat (\C, N \times N), \qquad B = B_0 + B_1,
$$
where $B_0$ is a diagonal part of $B$ and $B_1$ is an off-diagonal part of $B$.
Assume that
\begin{equation}
\label{318}
B_0 \ge a I_{N \times N} \qquad \text{and} \qquad
\|B_1\|_{S_2}^2 \le b (N^2-N),
\end{equation}
with some positive $a$ and $b$.
Then
$$
N_+ (B) \ge N - \frac{b(N^2-N)}{a^2}.
$$
\end{lemma}

\begin{proof}
The second condition \eqref{318} implies that
$$
\# \left(\spec (B_1) \cap (-\infty, -a]\right) \le \frac{b(N^2-N)}{a^2}.
$$
Therefore,
$$
\# \left(\spec (B) \cap (0, \infty)\right)
\ge N - \frac{b(N^2-N)}{a^2}.
\quad \qedhere
$$
\end{proof}

\begin{lemma}
\label{l319}
Fix $a, b > 0$. 
Assume that there is a sequence of matrices $B^{(N)}$ 
satisfying the conditions of Lemma \ref{l318} for each $N \in \N$.
Then one can find such a number $N$ that
$$
N_+ (B^{(N)}) \ge \frac{a^2}{4b} + \frac12.
$$
\end{lemma}

\begin{proof}
We take as $N$ a nearest integer to the number $\frac{a^2 + b}{2b}$,
so
$$
N = \frac{a^2 + b}{2b} + \er, \qquad |\er| \le \frac12.
$$
Then we use Lemma \ref{l318}, and we have
\begin{eqnarray*}
N - \frac{b(N^2-N)}{a^2} = \frac{(a^2+b)N}{a^2} - \frac{bN^2}{a^2}
= \frac{a^2+b}{a^2} \left(\frac{a^2 + b}{2b} + \er\right) 
- \frac{b}{a^2} \left(\frac{a^2 + b}{2b} + \er\right)^2 \\
= \frac{(a^2+b)^2}{4a^2b} - \frac{b\er^2}{a^2}
\ge \frac{a^4 + 2a^2b + b^2}{4a^2b} - \frac{b}{4a^2} 
= \frac{a^2}{4b} + \frac12. \quad \qedhere
\end{eqnarray*}
\end{proof}

\subsection{Proof of Theorem \ref{t14}}
Recall that for each natural $N$ Lemma \ref{l317} guarantees the existence of a set 
$H^{(N)} \in \left(S_{\sqrt{\la-\de^2}}^{d-2}\right)^N$
such that the corresponding space $F(H^{(N)})$ possesses the following properties:
\begin{itemize}
\item $F(H^{(N)}) \subset W_2^1 (\O)$;
\item $F (H) \cap H_{\cal D} (\la) = \{0\}$;
\item $\dim F(H^{(N)}) = N$;
\item $- \D f = \la f$ for all $f \in F(H^{(N)})$;
\item if 
$f = \sum_{j=1}^N c_j u_{\eta_j} \in F(H^{(N)})$ then
\begin{equation}
\label{319}
\la \int_\O |f(x)|^2 dx - \int_\O |\n f(x)|^2 dx = (B^{(N)} \vec c, \vec c),
\end{equation}
with $B^{(N)} = B_0^{(N)} + B_1^{(N)}$, 
$B_0^{(N)}$ satisfies \eqref{38},
$B_1^{(N)}$ satisfies \eqref{317}.
\end{itemize}
We apply Lemma \ref{l319} with
$$
a = \frac{\mes_d \O}{2 w(\O)^2}, \quad
b = \frac{32 (d-4)^{\frac{d-4}2} \oo_{d-2} (2\pi)^{d-2} \la^{\frac{d}2}}
{(d-1) d^{\frac{d}2} \oo_{d-1}^2 (\la-\de^2)^{d-3}}  
\left(\oo_{d-1} w(\O)^{d-2} \left(\mes_d\O\right)^d\right)^{\frac1{d-1}} .
$$
It means that we can find a number $N$, the set 
$H^{(N)} \in \left(S_{\sqrt{\la-\de^2}}^{d-2}\right)^N$,
and a subspace $F_* \subset F(H^{(N)})$ such that
the quadratic form \eqref{319} is positive on this subspace $F_*$, and
\begin{eqnarray*}
\dim F_* > \frac{a^2}{4b} 
= \frac{(\mes_d \O)^2}{16 w(\O)^4} \cdot
\frac{(d-1) d^{\frac{d}2} \oo_{d-1}^2 (\la-\de^2)^{d-3}}
{32 (d-4)^{\frac{d-4}2} \oo_{d-2} (2\pi)^{d-2} \la^{\frac{d}2}
\left(\oo_{d-1} w(\O)^{d-2} \left(\mes_d\O\right)^d\right)^{\frac1{d-1}}} \\
\ge \frac{d^3 (\la-\de^2)^{d-3} \oo_{d-1}^{2 - \frac1{d-1}} (\mes_d \O)^{\frac{d-2}{d-1}}}
{512 \la^{\frac{d}2} (2\pi)^{d-2} \oo_{d-2} w(\O)^{5 - \frac1{d-1}}}
= \frac{d^3 (\la-\de^2)^{d-3} \G\left(\frac{d}2\right) (\mes_d \O)^{\frac{d-2}{d-1}}}
{2^{d+7} \pi^{\frac{d-3}2} \la^{\frac{d}2} 
\G \left(\frac{d+1}2\right)^{2 - \frac1{d-1}} w(\O)^{5 - \frac1{d-1}}}.
\end{eqnarray*}
This estimate together with Lemma \ref{l31} gives

\begin{lemma}
\label{l320}
Let $d\ge 5$. 
Assume that a bounded domain $\O$ is contained in a layer of width $w(\O)$.
Assume that $\la \ge \de^2$ where $\de = w(\O)^{-1}$.
Then
\begin{equation}
\label{320}
N_{\cal N} (\O, \la) - N_{\cal D} (\O, \la) \ge
\frac{d^3 (\la-\de^2)^{d-3} \G\left(\frac{d}2\right) (\mes_d \O)^{\frac{d-2}{d-1}}}
{2^{d+7} \pi^{\frac{d-3}2} \la^{\frac{d}2} 
\G \left(\frac{d+1}2\right)^{2 - \frac1{d-1}} w(\O)^{5 - \frac1{d-1}}}.
\end{equation}
\end{lemma}
 
{\it Proof of Theorems \ref{t335} and \ref{t14}.}
By virtue of \eqref{15} it is sufficient to estimate the right hand side of \eqref{320}
from below for $\la \ge \la_1 (\O)$.
By the Faber-Krahn inequality
$$
\la_1 (\O) \ge \frac{j_{\frac{d}2-1}^2 \oo_d^{2/d}}{(\mes_d \O)^{2/d}},
$$
thus, in the rest of the proof we assume
\begin{equation}
\label{321}
\la \ge \frac{j_{\frac{d}2-1}^2 \oo_d^{2/d}}{(\mes_d \O)^{2/d}}.
\end{equation}

Let us find the asymptotics of the right hand side of \eqref{320}
if $d \to \infty$ up to terms of order $O (\log d)$.
Recall that 
$$
\de = \frac1{w(\O)} = \frac1{\left(d! \mes_d \O\right)^{1/d}},
$$
so,
$$
\frac{\de^2}\la \le \frac{(\mes_d \O)^{2/d}}
{\left(d! \mes_d \O\right)^{2/d} j_{\frac{d}2-1}^2 \oo_d^{2/d}}
= \frac{\G\left(\frac{d}2 + 1\right)^{2/d}}{\pi (d!)^{\frac2d} \,j_{\frac{d}2-1}^2}
= O \left(\frac1{d^3}\right), \quad d\to \infty.
$$
Therefore,
$$
\log \left(\frac{\la-\de^2}\la\right) = O \left(\frac1{d^3}\right),
$$
and
\begin{eqnarray}
\nonumber
\log \left(\frac{(\la-\de^2)^{d-3}}{\la^{\frac{d}2}}\right) 
= \left(\frac{d}2 - 3\right) \log \la + (d-3) \log \left(\frac{\la-\de^2}\la\right)\\
\label{322}
= \left(\frac{d}2 - 3\right) \log \la + O \left(\frac1{d^2}\right) \\
\ge \left(\frac{d}2 - 3\right)
\left(2 \log j_{\frac{d}2-1} + \frac2d \log \oo_d - \frac2d \log \mes_d \O\right)
+ O \left(\frac1{d^2}\right), \quad d\to \infty,
\nonumber
\end{eqnarray}
where we used \eqref{321}.

We have also
$$
\log \oo_d = \frac{d}2 \log \pi - \log \G \left(\frac{d}2 + 1\right)
= \frac{d}2 \log \pi - \frac{d}2 \left(\log d - \log 2 - 1\right) + O (\log d), \quad d\to \infty.
$$
Substituting this and \eqref{24} in \eqref{322} we obtain
\begin{eqnarray}
\nonumber
\log \left(\frac{(\la-\de^2)^{d-3}}{\la^{\frac{d}2}}\right) 
\ge (d-6) \left(\log j_{\frac{d}2-1} + \frac1d \log \oo_d - \frac1d \log \mes_d \O\right)
+ O \left(\frac1{d^2}\right) \\
\label{323}
= \frac{d \log d}2 + \frac{d}2 \left(\log \pi - \log 2 +1\right)
+ |a_1| 2^{1/3} d^{1/3} - \frac{d-6}d \, \log \mes_d \O + O (\log d), \\
d\to \infty.
\nonumber
\end{eqnarray}
Next,
\begin{equation}
\label{324}
\log \G \left(\frac{d}2\right) 
= \frac{d \log d}2 - \frac{d}2 \left(\log 2 +1\right) + O (\log d), 
\end{equation}
\begin{eqnarray}
\nonumber
\log\left(\G \left(\frac{d+1}2\right)^{2 - \frac1{d-1}}\right)
= \left(2 - \frac1{d-1}\right) \log \G \left(\frac{d+1}2\right) \\
\label{325}
= d \log d - d \left(\log 2 + 1\right) + O (\log d), \quad d \to \infty.
\end{eqnarray}
Finally, 
\begin{eqnarray}
\nonumber
\log \left(\frac{(\mes_d \O)^{\frac{d-2}{d-1}}}{w(\O)^{5 - \frac1{d-1}}}\right)
= \frac{d-2}{d-1} \log \mes_d \O 
- \frac{5d-6}{d(d-1)} \log \left(d! \mes_d \O\right) \\
= \frac{d-6}d \, \log \mes_d \O + O (\log d), \quad d \to \infty.
\label{326}
\end{eqnarray}
Substituting \eqref{323}, \eqref{324}, \eqref{325} and \eqref{326} into \eqref{320} we get
\begin{eqnarray*}
\log \Psi (d, k, \O) 
\ge  \frac{d \log d}2 + \frac{d}2 \left(\log \pi - \log 2 +1\right)
+ |a_1| 2^{1/3} d^{1/3} - \frac{d-6}d \, \log \mes_d \O \\
+ \frac{d \log d}2 - \frac{d}2 \left(\log 2 +1\right) 
- d \log 2 - \frac{d}2 \log \pi \\
- d \log d + d \left(\log 2 + 1\right) 
+ \frac{d-6}d \, \log \mes_d \O + O (\log d) \\
= \left(1 - \log 2\right) d + |a_1| 2^{1/3} d^{1/3} + O (\log d), \quad d \to \infty.
\end{eqnarray*}
Therefore,
$$
\log \Psi (d, k, \O) \ge \left(1 - \log 2\right) d
$$
for sufficiently large $d$.
Taking into account \eqref{*} we see that there is a constant $C_{conv} > 0$ such that
$$
\Psi (d,k,\O) \ge C_{conv} \left(\frac{e}2\right)^d \quad \text{for all}\ k\ \text{and}\ \O.
\qquad \qed
$$

%%%%%%%%%%%%%%%%%%%%%%%%%%%%%%%%%%%%%%%%%%
\section{Ball}
In this section we prove Theorem \ref{t16}.
We use some ideas from \cite{CMS}.

Let $\O = B_1 (0) \subset \R^d$ be a unit ball centered at the origine.
The first Dirichlet eigenvalue is 
$\la_1 = j_{\frac{d}2-1}^2$, see \eqref{21}.
The first eigenfunction of the Neumann problem is $\psi_1 (x) \equiv 1$
with the eigenvalue $\mu_1 = 0$.
Other eigenfunctions are
$$
\psi_{m,k} (x) = r^{1-\frac{d}2} J_{\frac{d}2 + m -1} \left(p_{d,m,k} r\right) Y_m (\oo),
\quad m \in \N_0, \ k \in \N.
$$
Here $x = (r; \oo)$ are spherical coordinates in $\R^d$,
$Y_m$ are sphercial harmonics,
and $p_{d,m,k}$ is the $k$-th positive root of the function
\begin{equation}
\label{41}
\left(r^{1-\frac{d}2} J_{\frac{d}2 + m -1} (r)\right)'
= r^{-\frac{d}2} \left(r J_{\frac{d}2 + m -1}' (r) 
+ \left(1-\frac{d}2\right) J_{\frac{d}2 + m -1} (r)\right).
\end{equation}
The corresponding eigenvalues are 
$$
\mu_{m,k} = \left(p_{d,m,k}\right)^2
\quad \text{with multiplicity} \quad 
\ka_m := \left( \begin{array}{cc} m + d - 1 \\ d-1 \end{array} \right) 
- \left( \begin{array}{cc} m + d - 3 \\ d-1 \end{array} \right) ,
$$
where we assume $\left( \begin{array}{cc} p \\ q \end{array} \right) = 0$
if $p < q$.

By virtue of Dixon's theorem \cite[\S 15.23]{Watson} the roots of the functions
$$
r J_{\frac{d}2 + m -1}' (r) + \left(1-\frac{d}2\right) J_{\frac{d}2 + m -1} (r)
\quad \text{and} \quad J_{\frac{d}2 + m -1} (r)
$$
interlace.
Thus, for all natural $m$ we have $p_{d,m,2} > j_{\frac{d}2+m-1}$
due to \eqref{41}.
Next, $j_{\nu+1} > j_\nu$, see for example \cite[\S 15.22]{Watson},
and therefore, $p_{d,m,2} > j_{\frac{d}2-1}$.
So, the set of eigenvalues $\mu_{m,k}$ that lesser than $\la_1$
consists of the eigenvalues of type $\mu_{m,1} = p_{d,m,1}^2$ only.
Note also that for $m=0$
$$
\left(r^{1-d/2} J_{\frac{d}2-1} (r)\right)' = - r^{-d/2} J_{\frac{d}2} (r),
$$
so, $p_{d,0,1} = j_{\frac{d}2} > j_{\frac{d}2-1}$.
The numbers $p_{d,m,1}$ increase in $m$ if $m \ge 1$, see \eqref{ratio} below.
Denote 
\begin{equation}
\label{42}
M = M(d) = \max \left\{m \in \N : p_{d,m,1} \le j_{\frac{d}2-1}\right\}.
\end{equation}
Taking into account the first Neumann eigenvalue $\mu_1=0$ we get
$$
\Phi (d,1,B) = 1 + \sum_{m=1}^M \ka_m = \sum_{m=0}^M \ka_m
= \left( \begin{array}{cc} M + d - 1 \\ d-1 \end{array} \right)  
+ \left( \begin{array}{cc} M + d - 2 \\ d-1 \end{array} \right) .
$$
This equality implies
$$
\left( \begin{array}{cc} M + d - 1 \\ d-1 \end{array} \right) 
< \Phi (d,1,B) < 2 \left( \begin{array}{cc} M + d - 1 \\ d-1 \end{array} \right) ,
$$
and
\begin{equation}
\label{43}
\log \Phi (d,1,B) = \log \left( \begin{array}{cc} M + d - 1 \\ d-1 \end{array} \right) 
+ O (1), \quad d \to \infty.
\end{equation}

In order to find the asymptotics of $M(d)$ for large $d$ we need 
to describe the behaviour of $p_{d,m,1}$.
If $m \ge 1$ the numbers $p_{d,m,1}^2$ coincide with the minima
of the following ratio of quadratic forms
\begin{equation}
\label{ratio}
p_{d,m,1}^2 = \min \frac{\int_0^1\left(|f'(x)|^2 x^{d-1} 
+ m (m+d-2) |f(x)|^2 x^{d-3}\right) dx}
{\int_0^1|f(x)|^2 x^{d-1} dx},
\end{equation}
defined on the set
\begin{equation*}
\label{dom}
\dom = \left\{ f : 
\int_0^1\left(|f'(x)|^2 x^{d-1} + m (m+d-2) |f(x)|^2 x^{d-3} \right) < \infty \right\}.
\end{equation*}

\begin{lemma}
If $m\ge 1$ and $d\ge 3$ then
\begin{equation}
\label{45}
m(m+d-2) \le p_{d,m,1}^2 \le \frac{dm(m+d-2)}{d-2}.
\end{equation}
\end{lemma}

\begin{proof}
Clearly, the second term in the numerator of \eqref{ratio} satisfies
$$
\int_0^1 |f(x)|^2 x^{d-3} dx \ge \int_0^1|f(x)|^2 x^{d-1} dx ,
$$
which implies the first inequality \eqref{45}.

Next, we substitute $f(x) \equiv 1$ in \eqref{ratio}.
We get
$$
\int_0^1\left(|f'(x)|^2 x^{d-1} + m (m+d-2) |f(x)|^2 x^{d-3}\right) dx
%= m (m+d-2) \int_0^1 x^{d-3} dx 
= \frac{m(m+d-2)}{d-2},
$$
$$
\int_0^1|f(x)|^2 x^{d-1} dx = \frac1d,
$$
which implies the claim.
\end{proof}

\begin{rem}
The lower bound \eqref{45} is just the same as (1'') in \cite{LS}.
The upper bound \eqref{45} is better than (1) in \cite{LS} for $d \ge 3$.
\end{rem}

The asymptotics \eqref{23} yields
$$
\la_1 = j_{\frac{d}2-1}^2 = \frac{d^2}4 + |a_1| 2^{-2/3} d^{4/3} + O (d),
\quad d \to \infty.
$$
The equations
$$
m (m+d-2) = \la_1 \quad \text{and} \quad
m (m+d-2) = \frac{(d-2)\la_1}{d}
$$
have positive solutions
$$
m^* = \frac12 \left(\sqrt{4\la_1 + (d-2)^2} - d + 2\right)
\quad \text{and} \quad
m_* = \frac12 \left(\sqrt{\frac{4\la_1(d-2)}{d} + (d-2)^2} - d + 2\right)
$$
respectively.
Therefore, the number $M(d)$ defined by \eqref{42} satisfies the bounds
$$
m_*-1 < M (d) \le m^*.
$$
Furthermore,
$$
m^* = \frac12 \left(\sqrt{2 d^2 + |a_1| 2^{4/3} d^{4/3} + O (d)} - d + 2\right)
= \frac{\sqrt 2 - 1}2 \, d + |a_1| 2^{-7/6} d^{1/3} + O (1), \quad d \to \infty.
$$
In the same way,
$$
m_* = \frac{\sqrt 2 - 1}2 \, d + |a_1| 2^{-7/6} d^{1/3} + O (1), \quad d \to \infty,
$$
and thus, the same is true for $M$,
$$
M (d) = \frac{\sqrt 2 - 1}2 \, d + |a_1| 2^{-7/6} d^{1/3} + O (1), \quad d \to \infty.
$$
Clearly,
\begin{equation}
\label{46}
\log \left( \begin{array}{cc} M + d - 1 \\ d-1 \end{array} \right) 
= \log \left((M+d-1)!\right) - \log \left(M!\right) - \log\left((d-1)!\right).
\end{equation}

The Stirling formula implies the following

\begin{lemma}
\label{l42}
If 
$$
K(d) = \ga_0 d + \ga_1 d^{1/3} + O(1), \quad d \to \infty, \quad \ga_0 > 0,
$$
then
$$
\log \left(K(d)!\right) = \ga_0 d \log d + \ga_0 (\log \ga_0 -1) d 
+ \ga_1 d^{1/3} \log d + \ga_1 \log \ga_0 d^{1/3} + O (\log d), \quad d \to \infty.
$$
\end{lemma}

Therefore,
\begin{eqnarray}
\nonumber
\log \left((M+d-1)!\right) 
= \frac{\sqrt 2 + 1}2\, d \log d + \frac{\sqrt 2 + 1}2 \left(\log \frac{\sqrt 2 + 1}2 - 1\right) d \\
+ |a_1| 2^{-7/6} d^{1/3} \log d + |a_1| 2^{-7/6} \log \frac{\sqrt 2 + 1}2 d^{1/3} + O(\log d),
\label{47}
\end{eqnarray}
\begin{eqnarray}
\nonumber
\log \left(M!\right) 
= \frac{\sqrt 2 - 1}2\, d \log d + \frac{\sqrt 2 - 1}2 \left(\log \frac{\sqrt 2 - 1}2 - 1\right) d \\
+ |a_1| 2^{-7/6} d^{1/3} \log d + |a_1| 2^{-7/6} \log \frac{\sqrt 2 - 1}2 d^{1/3} + O(\log d),
\label{48}
\end{eqnarray}
and
\begin{equation}
\label{49}
\log \left((d-1)!\right)
= d \log d - d + O (\log d), \quad d \to \infty.
\end{equation}
Substituting \eqref{46}, \eqref{47}, \eqref{48}, \eqref{49} into \eqref{43} we get
\begin{equation}
\label{410}
\log \Phi (d,1,B) 
= \left(\sqrt 2 \log (\sqrt 2 + 1) - \log 2\right) d
+ |a_1| 2^{-1/6} \log (\sqrt 2 + 1) d^{1/3} + O (\log d), \quad d \to \infty,
\end{equation}
where $a_1$ is the first zero of the Airy function.
Theorem \ref{t16} is proved.
\qed

%%%%%%%%%%%%%%%%%%%%%%%%%%%%%%%%%%%%%%%%%%
\begin{thebibliography}{99}

\bibitem{Av86} P.~Aviles, 
{\it Symmetry theorems related to Pompeiu’s problem}, 
Amer. J. Math. 108 (1986), no. 5, 1023--1036.

\bibitem{Bezdek} K.~Bezdek, 
{\it Tarski's plank problem revisited}, 
Geometry --- intuitive, discrete, and convex, 45--64.
Bolyai Soc. Math. Stud., 24,
János Bolyai Math. Society, Budapest, 2013.

\bibitem{CMS}
G.~Cox, S.~MacLachlan, L.~Steeves, 
{\it Isoperimetric relations between Dirichlet and Neumann eigenvalues}, 
arXiv:1906.10061 (2019).

\bibitem{Faber} G.~Faber, 
{\it Beweis, dass unter allen homogenen Membranen von gleicher Fläche 
und gleicher Spannung die kreisförmige den tiefsten Grundton gibt}, 
Sitzungberichte der mathematisch-physikalischen Klasse 
der Bayerischen Akademie der Wissenschaften zu München (1923), 169--172.

\bibitem{F04} N.~Filonov,
{\it On an inequality for the eigenvalues of the Dirichlet
and Neumann problems for the Laplace operator}, 
Algebra i Analiz 16 (2004), no. 2, 172--176 (Russian). 
Translation in St. Petersburg Math. J. 16 (2005), no. 2, 413--416.

\bibitem{FilSaf} N.~Filonov, Y.~Safarov, 
{\it Asymptotic estimates for the difference between the Dirichlet and Neumann counting functions}, Funktsional. Anal. i Prilozhen., 44 (2010), no. 4, 54--64 (Russian); 
Engl. translation in Funct. Anal. Appl., 44 (2010), no. 4, 286--294.

\bibitem{Firey} W.~J.~Firey, 
{\it Lower bounds for volumes of convex bodies}, 
Arch. Math. 16 (1965), 69--74.

\bibitem{FrG} P.~Freitas, M.~Gama,
{\it On the (growing) gap between Dirichlet and Neumann eigenvalues},
arXiv:2405.18079 (2024).

\bibitem{Fr91} L.~Friedlander, 
{\it Some inequalities between Dirichlet and Neumann eigenvalues},
 Arch. Rational Mech. Anal. 116 (1991), no. 2, 153--160.
 
\bibitem{Hatcher} L.~Hatcher,
{\it Geometric inequalities between Dirichlet and Neumann eigenvalues},
arXiv:2504.18517 (2025).

\bibitem{Krahn}
E.~Krahn,
{\it Über eine von Rayleigh formulierte Minimaleigenschaft des Kreises}, 
Math. Ann. 94 (1925), no. 1, 97--100. 

\bibitem{Kra} 
E.~Krahn,
{\it Über Minimaleigenschaften der Kugel in drei un mehr Dimensionen},
Acta Comm. Univ. Dorpat. A9 (1926), 1--44.

\bibitem{Kr} P.~Kr\" oger,
{\it Upper bounds for the Neumann eigenvalues on a bounded domains in Euclidean space}, 
J. Funct. Anal. 106 (1992), 353--357.

\bibitem{LW86} H.~A.~Levine, H.~F.~Weinberger, 
{\it Inequalities between Dirichlet and Neumann eigenvalues}, 
Arch. Rational Mech. Anal. 94 (1986), 193--208.

\bibitem{LS}
L.~Lorch, P.~Szego, 
{\it Bounds and monotonicities for the zeros of derivatives of ultraspherical Bessel functions},
SIAM J.Math. Anal. 25 (1994), no. 2, 549--554.

\bibitem{Pal} J.~P\'al, 
{\it Ein minimumproblem f\" ur Ovale}, 
Math. Ann. 83 (1921), 311--319.

\bibitem{Payne55} L.~E.~Payne, 
{\it Inequalities for eigenvalues of membranes and plates}, 
J. Rational Mech. Anal. 4 (1955), 517--529.

\bibitem{P52} G.~P\'olya, 
{\it Remarks on the foregoing paper}, 
J. Math. Physics 31 (1952), 55--57.

\bibitem{QuWong}
C.~K.~Qu, R.~Wong,
{\it “Best possible” upper and lower bounds for the zeros of 
the Bessel function $J_\nu(x)$},
Trans. Amer. Math. Soc. 351 (1999), no 7, 2833--2859.

\bibitem{R25} J.~Rohleder, 
{\it Inequalities between Neumann and Dirichlet Laplacian eigenvalues on planar domains},
Math. Ann. 392 (2025), no. 4, 5553--5571.

\bibitem{Watson}
G.~N.~Watson, 
{\it A treatise on the theory of Bessel functions}, 
Cambridge University Press, Cambridge, 1995.




\end{thebibliography}
\end{document}
















